
\section{计算几何}
\subsection{极角排序}
\begin{lstlisting}[caption={}]
#include<bits/stdc++.h>
using i64 = long long;
const double eps = 0,pi = acos(-1);
struct node{
    i64 x,y,idx;
    double dge;
};
bool cmp(node a,node b){
    return a.dge - b.dge < eps;
}
int main(){
	std::ios::sync_with_stdio(0),std::cin.tie(0),std::cout.tie(0);
    i64 n;std::cin >> n;
    std::vector<node> a(n + 1);
    for(int i = 1;i <= n;i++){
        i64 x,y;std::cin >> x >> y;
        a[i] = {x,y,i,atan2(x,y)};
        if(a[i].dge < 0) a[i].dge += 2 * pi;
    }
    sort(a.begin() + 1,a.end(),cmp);
    double ans = a[1].dge - a[n].dge;
    ans += 2 * pi;
    i64 idx1 = a[1].idx,idx2 = a[n].idx;
    for(int i = 1;i < n;i++){
        double temp = a[i + 1].dge - a[i].dge;
        if(temp - ans < eps){
            ans = temp;
            idx1 = a[i + 1].idx,idx2 = a[i].idx;
        } 
    }
    std::cout << idx1 << " " << idx2 << "\n";
	return 0;
}
\end{lstlisting}

\subsection{二维静态凸包(Andrew 算法)}
$O(n \log n)$
\begin{lstlisting}[caption={}]
using Point = std::array<i64,2>;
using Vector = std::array<i64,2>;
Vector operator-(const Point& a,const Point& b){
	return {a[0] - b[0],a[1] - b[1]};
}
i64 operator*(const Vector& a,const Vector& b){
	return a[0] * b[1] - a[1] * b[0];
}
struct Convex{
	//逆时针
	std::vector<Point> points,convex;
	void addPoint(i64 x,i64 y){
		points.push_back({x,y});
	}
	void work(){
		std::sort(points.begin(),points.end());
		std::vector<int> stk;
		stk.push_back(0);
		for(int i = 1;i < points.size();i++){
			int m = stk.size();
			while(m >= 2 && (points[stk[m - 2]] - points[stk[m - 1]]) * (points[stk[m - 1]] - points[i]) <= 0){
				m--;
				stk.pop_back();
			}
			stk.push_back(i);
		}
		int temp = stk.size();
		for(int i = points.size() - 1;i >= 0;i--){
			int m = stk.size();
			while(m > temp && (points[stk[m - 2]] - points[stk[m - 1]]) * (points[stk[m - 1]] - points[i]) <= 0){
				m--;
				stk.pop_back();
			}
			stk.push_back(i);
		}
		for(int i = 0;i < stk.size() - 1;i++){
			if(i >= 1 && points[stk[i]] == points[stk[i - 1]]) continue;
			convex.push_back(points[stk[i]]);
		}
	}
};
\end{lstlisting}




\subsection{闵可夫斯基和}
计算两个凸包合成的大凸包。
\begin{lstlisting}[caption={}]
template<class T> vector<Point<T>> mincowski(vector<Point<T>> P1, vector<Point<T>> P2) {
    int n = P1.size(), m = P2.size();
    vector<Point<T>> V1(n), V2(m);
    for (int i = 0; i < n; i++) {
        V1[i] = P1[(i + 1) % n] - P1[i];
    }
    for (int i = 0; i < m; i++) {
        V2[i] = P2[(i + 1) % m] - P2[i];
    }
    vector<Point<T>> ans = {P1.front() + P2.front()};
    int t = 0, i = 0, j = 0;
    while (i < n && j < m) {
        Point<T> val = sign(cross(V1[i], V2[j])) > 0 ? V1[i++] : V2[j++];
        ans.push_back(ans.back() + val);
    }
    while (i < n) ans.push_back(ans.back() + V1[i++]);
    while (j < m) ans.push_back(ans.back() + V2[j++]);
    return ans;
}
\end{lstlisting}


\subsection{平面最近点对}
$O(n \log n)$
\begin{lstlisting}[caption={}]
struct CPOP{
    using Point = std::array<i64,2>;
    std::vector<Point> points;   
    i64 minDist = 9e18;
    Point p1,p2; 
    i64 dist(Point p1,Point p2){
        auto [x1,y1] = p1;
        auto [x2,y2] = p2;
        return (x1 - x2) * (x1 - x2) + (y1 - y2) * (y1 - y2);
    }
    void addPoint(i64 x,i64 y){
        points.push_back({x,y});
    }
    //0 idx 左闭右开
    void work(i64 l,i64 r){
        if(r - l == 1) return;
        int mid = (l + r) >> 1;
        i64 midx = points[mid][0];
        work(l,mid),work(mid,r);

        std::inplace_merge(points.begin() + l,points.begin() + mid,points.begin() + r,
                        [&](Point a,Point b){return a[1] < b[1];});
        
        std::vector<Point> temp;
        for(int i = l;i < r;i++){
            if((points[i][0] - midx) * (points[i][0] - midx) < minDist) 
                temp.push_back(points[i]);
        }

        for(int i = 0;i < temp.size();i++){
            for(int j = i + 1;j < temp.size();j++){
                if((temp[i][1] - temp[j][1]) * (temp[i][1] - temp[j][1]) >= minDist) break;
                i64 d = dist(temp[i],temp[j]);
                if(d < minDist){
                    minDist = d;
                    p1 = temp[i],p2 = temp[j];
                }
            }
        }
    }
    void work(){
        std::sort(points.begin(),points.end());
        work(0,points.size());
    }
};
\end{lstlisting}

\subsection{扫描线}
矩形面积并
\begin{lstlisting}[caption={}]
struct ScanLine{
	std::vector<i64> cover,weight,num;
	std::vector<std::array<i64,4>> side;
	std::map<i64,i64> mp;
	i64 n;
	ScanLine(){};
	ScanLine(i64 N):n(N),cover(N << 3),weight(N << 3),num((N << 1) + 1){}
    void push_up(i64 p,i64 l,i64 r){
        if(cover[p]) weight[p] = num[r] - num[l];
		else if(l + 1 == r) weight[p] = 0;
		else weight[p] = weight[(p << 1) | 1] + weight[p << 1];
    }
    void update(i64 l,i64 r,i64 c,i64 s,i64 t,i64 p){
        if(l <= s && t <= r){
            cover[p] += c;
			push_up(p,s,t);
            return;
        }
		if(s + 1 == t) return ;
        i64 m = s + ((t - s) >> 1);
        if(l <= m) update(l,r,c,s,m,p << 1);
        if(r > m) update(l,r,c,m,t,(p << 1) | 1);
        push_up(p,s,t);
    }
	void addRectangle(i64 x1,i64 y1,i64 x2,i64 y2){
		side.push_back({x1,y1,y2,1});
		side.push_back({x2,y1,y2,-1});
		mp[y1]++,mp[y2]++;
	}
	i64 run(){
		i64 res = 0,idx = 0;
		for(auto& i : mp){
			i.second = ++idx;
			num[idx] = i.first;
		}
		sort(side.begin(),side.end());
		update(mp[side[0][1]],mp[side[0][2]],1,1,2 * n,1);
		for(int i = 1;i < side.size();i++){
			auto [x,y1,y2,tag] = side[i];
			res += (x - side[i - 1][0]) * weight[1];
			update(mp[y1],mp[y2],tag,1,2 * n,1);
		}
		return res;
	}
};

int main(){
	std::ios::sync_with_stdio(0),std::cin.tie(0),std::cout.tie(0);
	i64 n;
	std::cin >> n;
	ScanLine sl(n);
	for(int i = 1;i <= n;i++){
		i64 x1,y1,x2,y2;
		std::cin >> x1 >> y1 >> x2 >> y2;
		sl.addRectangle(x1,y1,x2,y2);
	}
	std::cout << sl.run() << "\n";
	return 0;
}
\end{lstlisting}

\subsection{平面最远点对(旋转卡壳)}
$O(N)$,需要凸包
\begin{lstlisting}[caption={}]
using Point = std::array<i64,2>;
using Vector = std::array<i64,2>;
Vector operator-(const Point& a,const Point& b){
	return {a[0] - b[0],a[1] - b[1]};
}
i64 operator*(const Vector& a,const Vector& b){
	return a[0] * b[1] - a[1] * b[0];
}
//Furthest Pair of Points O(N)
i64 FPOP(std::vector<Point>& points){
	i64 j = 0,ans = 0,m = points.size();
	if(m == 1) return 0;
	else if(m == 2) return dist(points[0],points[1]);
	for(int i = 0;i < m;i++){ //需要枚举边而不是点
		int las = i == 0 ? m - 1 : i - 1;
		//通过叉乘算三角形面积来替代点到直线距离
		while(std::abs((points[las] - points[j]) * (points[las] - points[i])) <=
			std::abs((points[las] - points[(j + 1) % m]) * (points[las] - points[i]))) j = (j + 1) % m;
		ans = std::max({ans,dist(points[j],points[i]),dist(points[j],points[las])});
	}
	return ans;
}
\end{lstlisting}

\subsection{最小矩形覆盖(旋转卡壳)}
$O(N)$,需要凸包
\begin{lstlisting}[caption={}]
i64 MRC(std::vector<Point>& points){
	i64 j = 1,l = 0,r = 0,ans = 1e18,m = points.size();
	if(m <= 2) return 0;
	for(int i = 0;i < m;i++){ //需要枚举边而不是点
		int las = i == 0 ? m - 1 : i - 1;
		//通过叉乘算三角形面积来替代点到直线距离
		while(std::abs((points[las] - points[j]) * (points[las] - points[i])) <=
			std::abs((points[las] - points[(j + 1) % m]) * (points[las] - points[i]))) j = (j + 1) % m;
		//通过点乘投影来判断长度
		while(((points[i] - points[las]) ^ (points[r] - points[las])) <=
			((points[i] - points[las]) ^ (points[(r + 1) % m] - points[las]))) r = (r + 1) % m;
		if(i == 0) l = r;
		while(((points[las] - points[i]) ^ (points[l] - points[i])) <=
			((points[las] - points[i]) ^ (points[(l + 1) % m] - points[i]))) l = (l + 1) % m;

		i64 t1 = std::abs((points[las] - points[j]) * (points[i] - points[j]));
		i64 t2 = std::abs((points[r] - points[las]) ^ (points[i] - points[las])) + 
				 std::abs((points[l] - points[i]) ^ (points[las] - points[i]));
		i64 t3 = std::abs((points[las] - points[i]) ^ (points[las] - points[i]));
		ans = std::min<i64>(ans,(__int128)t1 * (t2 - t3) / t3);
	}
	return ans;
}
\end{lstlisting}

\subsection{最小圆覆盖}
期望$O(N)$,最坏$O(h^3 n)$
\begin{lstlisting}[caption={}]
#include<bits/stdc++.h>
using i64 = long long;
std::mt19937 rng(std::chrono::steady_clock::now().time_since_epoch().count());
static constexpr double eps = 1e-9;
using Point = std::array<double,2>;

double dist(Point& a,Point& b){
	return std::sqrt((a[0] - b[0]) * (a[0] - b[0]) + (a[1] - b[1]) * (a[1] - b[1]));
}

struct SEC{
	std::vector<Point> points;
	double r = 0;
	Point o{};
	void addPoint(double x,double y){
		points.push_back({x,y});
	}

	void update(Point& x,Point& y){
		o[0] = (x[0] + y[0]) / 2;
		o[1] = (x[1] + y[1]) / 2;
		r = dist(x,y) / 2;
	}

	Point geto(Point a,Point b,Point c){
		double a1,a2,b1,b2,c1,c2;
		Point res;
		a1 = 2 * (b[0] - a[0]), b1 = 2 * (b[1] - a[1]),
		c1 = b[0] * b[0] - a[0] * a[0] + b[1] * b[1] - a[1] * a[1];
		a2 = 2 * (c[0] - a[0]), b2 = 2 * (c[1] - a[1]),
		c2 = c[0] * c[0] - a[0] * a[0] + c[1] * c[1] - a[1] * a[1];
		if(std::fabs(a1) <= eps){
			res[1] = c1 / b1;
			res[0] = (c2 - res[1] * b2) / a2;
		}
		else if(std::fabs(b1) <= eps){
			res[0] = c1 / a1;
			res[1] = (c2 - res[0] * a2) / b2;
		}
		else{
			res[0] = (c2 * b1 - c1 * b2) / (a2 * b1 - a1 * b2);
			res[1] = (c2 * a1 - c1 * a2) / (b2 * a1 - b1 * a2);
		}
		return res;
	}

	void work(){
		int n = points.size();
		for(int i = 0;i < n;i++) std::swap(points[rng() % n],points[rng() % n]);
		o = points[0];
		for(int i = 0;i < n;i++){
			if(dist(o,points[i]) - r < eps) continue;
			update(points[i],points[0]);
			for(int j = 0;j < i;j++){
				if(dist(o,points[j]) - r < eps) continue;
				update(points[i],points[j]);
				for(int k = 0;k < j;k++){
					if(dist(o,points[k]) - r < eps) continue;
					o = geto(points[i],points[j],points[k]);
					r = dist(o,points[i]);	
				}
			}
		}
	}
};

int main(){
	std::ios::sync_with_stdio(0),std::cin.tie(0),std::cout.tie(0);
	int n;
	std::cin >> n;
	SEC circ;
	for(int i = 1;i <= n;i++){
		double x,y;
		std::cin >> x >> y;
		circ.addPoint(x,y);
	}
	circ.work();
	std::cout << std::fixed << std::setprecision(10) << circ.r << "\n" << circ.o[0] << " " << circ.o[1] << "\n";
	return 0;
}
\end{lstlisting}



\subsection{半平面交}
计算多条直线左边平面部分的交集。
\begin{lstlisting}[caption={}]
template<class T> vector<Point<T>> halfcut(vector<Line<T>> lines) {
    sort(lines.begin(), lines.end(), [&](auto l1, auto l2) {
        auto d1 = l1.b - l1.a;
        auto d2 = l2.b - l2.a;
        if (sign(d1) != sign(d2)) {
            return sign(d1) == 1;
        }
        return cross(d1, d2) > 0;
    });
    deque<Line<T>> ls;
    deque<Point<T>> ps;
    for (auto l : lines) {
        if (ls.empty()) {
            ls.push_back(l);
            continue;
        }
        while (!ps.empty() && !pointOnLineLeft(ps.back(), l)) {
            ps.pop_back();
            ls.pop_back();
        }
        while (!ps.empty() && !pointOnLineLeft(ps[0], l)) {
            ps.pop_front();
            ls.pop_front();
        }
        if (cross(l.b - l.a, ls.back().b - ls.back().a) == 0) {
            if (dot(l.b - l.a, ls.back().b - ls.back().a) > 0) {
                if (!pointOnLineLeft(ls.back().a, l)) {
                    assert(ls.size() == 1);
                    ls[0] = l;
                }
                continue;
            }
            return {};
        }
        ps.push_back(lineIntersection(ls.back(), l));
        ls.push_back(l);
    }
    while (!ps.empty() && !pointOnLineLeft(ps.back(), ls[0])) {
        ps.pop_back();
        ls.pop_back();
    }
    if (ls.size() <= 2) {
        return {};
    }
    ps.push_back(lineIntersection(ls[0], ls.back()));
    return vector(ps.begin(), ps.end());
}
\end{lstlisting}
